% ============================================================
%  Übungsaufgaben Template (my_dhbw Stil, uebung_*.tex)
%  Header "Übungsblatt", grün eingefärbte <details>-Lösungen.
%  Renderer: pdflatex (zweimal)
% ============================================================
\documentclass[11pt, a4paper]{article}

\usepackage[ngerman]{babel}
\usepackage{fontspec}
\setmainfont[
  Path=/usr/share/fonts/TTF/,
  Extension=.ttf,
  BoldFont=DejaVuSans-Bold,
  ItalicFont=DejaVuSans-Oblique,
  BoldItalicFont=DejaVuSans-BoldOblique
]{DejaVuSans}
\setsansfont[
  Path=/usr/share/fonts/TTF/,
  Extension=.ttf,
  BoldFont=DejaVuSans-Bold,
  ItalicFont=DejaVuSans-Oblique,
  BoldItalicFont=DejaVuSans-BoldOblique
]{DejaVuSans}
\setmonofont[Path=/usr/share/fonts/TTF/,Extension=.ttf]{DejaVuSansMono}
\usepackage[top=2cm, bottom=2cm, left=2.4cm, right=2.4cm, footskip=1cm]{geometry}
\usepackage{amsmath, amssymb}
\usepackage{xcolor}
\usepackage{enumitem}
\usepackage{fancyhdr}
\usepackage{lastpage}
\usepackage{tabularx}
\usepackage{booktabs}
\usepackage{microtype}
\usepackage{parskip}

\usepackage{array}
\usepackage{longtable}
\usepackage{ltablex}
\keepXColumns
\usepackage{colortbl}
\usepackage[most,breakable]{tcolorbox}
\usepackage{listings}
\usepackage[normalem]{ulem}
\usepackage{pdflscape}
\usepackage{titlesec}
\usepackage[colorlinks=true, linkcolor=accent, urlcolor=accent]{hyperref}

\renewcommand{\familydefault}{\sfdefault}

% ── Theme ─────────────────────────────────────────────────────
\definecolor{accent}{RGB}{0, 60, 113}
\definecolor{primaryblue}{RGB}{0, 60, 113}
\definecolor{loesung}{RGB}{0, 100, 0}
\definecolor{codebg}{RGB}{245, 245, 245}
\definecolor{figurebg}{RGB}{232, 241, 255}
\definecolor{tablehintbg}{RGB}{240, 255, 240}
\definecolor{solutionbg}{RGB}{240, 252, 240}
\definecolor{rulegray}{RGB}{180, 180, 180}
\definecolor{tableheadbg}{RGB}{0, 60, 113}

% ── Header/Footer ─────────────────────────────────────────────
\pagestyle{fancy}
\fancyhf{}
\renewcommand{\headrulewidth}{0.4pt}
\fancyhead[L]{\footnotesize\color{accent}\nichttitle}
\fancyhead[R]{\footnotesize\color{accent}\nichtsubtitle}
\fancyfoot[R]{\footnotesize\color{accent}Blatt \thepage\,/\,\pageref{LastPage}}

% ── Typografie ────────────────────────────────────────────────
\titleformat{\section}
    {\color{accent}\large\bfseries}{}{0pt}{}
    [\vspace{-4pt}\color{accent}\rule{\linewidth}{1.5pt}]
\titleformat{\subsection}
    {\normalsize\bfseries\color{accent!85!black}}{}{0pt}{}{}
\titleformat{\subsubsection}
    {\small\bfseries\color{accent!70!black}}{}{0pt}{}{}
\titlespacing*{\section}{0pt}{10pt}{3pt}
\titlespacing*{\subsection}{0pt}{6pt}{2pt}
\titlespacing*{\subsubsection}{0pt}{4pt}{1pt}

\setlength{\parindent}{0pt}
\setlength{\parskip}{6pt}
\renewcommand{\arraystretch}{1.2}
\setlist[itemize]{noitemsep, topsep=3pt, parsep=0pt, leftmargin=1.4em}
\setlist[enumerate]{noitemsep, topsep=3pt, parsep=0pt, leftmargin=1.6em}

% ── Boxen: solutionbox = Lösungsbox in grün ──────────────────
\newtcolorbox{figurebox}[1]{
  colback=figurebg, colframe=accent!50,
  title={Abbildung: #1}, fonttitle=\small\bfseries,
  breakable, before skip=6pt, after skip=6pt
}
\newtcolorbox{tablehintbox}[1]{
  colback=tablehintbg, colframe=green!50!black!50,
  title={Tabelle: #1}, fonttitle=\small\bfseries,
  breakable, before skip=6pt, after skip=6pt
}
\newtcolorbox{solutionbox}[1]{
  enhanced,
  colback=solutionbg, colframe=loesung,
  title={#1}, fonttitle=\small\bfseries,
  coltitle=white, colbacktitle=loesung,
  breakable, before skip=8pt, after skip=8pt
}

\lstset{
  basicstyle=\ttfamily\small,
  backgroundcolor=\color{codebg},
  breaklines=true,
  frame=single, framerule=0pt,
  xleftmargin=4pt, xrightmargin=4pt,
  aboveskip=6pt, belowskip=6pt,
  keepspaces=true
}

\providecommand{\nichttitle}{Übungsblatt}
\providecommand{\nichtsubtitle}{}

\renewcommand{\nichttitle}{Optimierungsverfahren}
\renewcommand{\nichtsubtitle}{Übungsblatt 5}


\begin{document}

\begin{center}
{\Large\bfseries\color{accent}\nichttitle}
\ifx\nichtsubtitle\empty\else
  \\[2pt]{\normalsize\color{accent!80!black}\nichtsubtitle}
\fi
\\[2pt]
\color{accent}\rule{0.4\linewidth}{0.8pt}\color{black}
\end{center}
\vspace{4pt}

% ── BODY ──────────────────────────────────────────────────────

\section{Übungsblatt 5 — 4. Nicht-Lineare Optimierung}


\noindent\textcolor{rulegray}{\rule{\linewidth}{0.4pt}}


\paragraph{Aufgabe 1 — Q-Matrix und Konvexität \textit{(10 Punkte)}}\mbox{}\\[2pt]

Gegeben sei das QO-Problem


\[
\min \tfrac{1}{2}\,\vec{x}^{\,T} Q\, \vec{x} + \vec{c}^{\,T}\vec{x}, \qquad Q = \begin{pmatrix} 4 & 2 \\ 2 & 3 \end{pmatrix},\ \vec{c} = \begin{pmatrix} -1 \\ -2 \end{pmatrix}
\]


ohne Nebenbedingungen.


a) Nenne die drei laut Kapitel relevanten Eigenschaften der Matrix $Q$ in einem QO-Problem und ihre jeweilige Konsequenz für die Lösbarkeit.

b) Prüfe für die obige Matrix $Q$, welche dieser Eigenschaften erfüllt sind. Begründe rechnerisch.

c) Welche Konsequenz folgt aus deiner Antwort in b) für das Optimierungsproblem? Bestimme zusätzlich die optimale Lösung $\vec{x}^*$ analytisch.


\textbf{Lösung:}


a)

\begin{itemize}
  \item \textbf{Symmetrie} ($Q = Q^T$): erforderlich, sonst lässt sich die Form $\tfrac{1}{2}x^TQx$ nicht eindeutig schreiben (man kann sie aber durch $\tfrac{1}{2}(Q+Q^T)$ symmetrisieren).
  \item \textbf{Positiv semi-definit} (keine negativen Eigenwerte): Zielfunktion ist konvex, jeder lokale Optimum ist global.
  \item \textbf{Positiv definit}: Zielfunktion strikt konvex → es existiert eine \textbf{eindeutige} globale Lösung.
\end{itemize}

b)

\begin{itemize}
  \item Symmetrie: $Q_{12} = Q_{21} = 2$ ✓
  \item Hauptminoren-Kriterium (Sylvester): $\det(4) = 4 > 0$, $\det(Q) = 4\cdot 3 - 2\cdot 2 = 8 > 0$ ✓
  \item Eigenwerte als Gegenprobe: $\lambda^2 - 7\lambda + 8 = 0 \Rightarrow \lambda_{1,2} = \tfrac{7 \pm \sqrt{17}}{2} \approx 5{,}56\ /\ 1{,}44$ — beide positiv.
\end{itemize}

$Q$ ist \textbf{symmetrisch und positiv definit}.


c) Damit ist das Problem strikt konvex und besitzt eine \textbf{eindeutige globale Lösung}. Da keine NB vorliegen, gilt $\nabla f = Q\vec{x} + \vec{c} = \vec{0}$:


\[
\vec{x}^* = -Q^{-1}\vec{c} = -\frac{1}{8}\begin{pmatrix} 3 & -2 \\ -2 & 4 \end{pmatrix}\begin{pmatrix} -1 \\ -2 \end{pmatrix} = -\frac{1}{8}\begin{pmatrix} 1 \\ -6 \end{pmatrix} = \begin{pmatrix} -1/8 \\ 3/4 \end{pmatrix}
\]



\noindent\textcolor{rulegray}{\rule{\linewidth}{0.4pt}}


\paragraph{Aufgabe 2 — CLS als QO formulieren \textit{(15 Punkte)}}\mbox{}\\[2pt]

Du sollst die Funktion $\hat{y}(x) = b_1 + b_2 x$ an folgende Daten anpassen:



{\small
\begin{tabularx}{\linewidth}{XXX}
\hline
\rowcolor{tableheadbg}
{\color{white}\textbf{$i$}} & {\color{white}\textbf{$x_i$}} & {\color{white}\textbf{$y_i$}} \\[2pt]
\hline
\endhead
\hline
\endfoot
1 & 1 & 2 \\
2 & 2 & 3 \\
3 & 3 & 5 \\
\end{tabularx}
}


Aus dem Anwendungskontext gilt: $b_1 \ge 0$, $b_2 \ge 0$, und es ist bekannt, dass $b_1 + b_2 = 4$.


a) Stelle Designmatrix $X$ und Beobachtungsvektor $\vec{y}$ auf und berechne $X^TX$ sowie $X^T\vec{y}$.

b) Warum kann hier nicht die analytische Lösung $\vec{b} = (X^TX)^{-1}X^T\vec{y}$ direkt verwendet werden? Begründe mit dem Kapitel.

c) Schreibe das Ausgleichsproblem als QO-Problem in Standardform $\min \tfrac{1}{2}\vec{b}^TQ\vec{b} + \vec{c}^T\vec{b}$ mit allen NB an. Gib $Q$, $\vec{c}$ und alle Restriktionen explizit an.


\textbf{Lösung:}


a) Mit Spalte für den Achsenabschnitt zuerst:


\[
X = \begin{pmatrix} 1 & 1 \\ 1 & 2 \\ 1 & 3 \end{pmatrix},\qquad \vec{y} = \begin{pmatrix} 2 \\ 3 \\ 5 \end{pmatrix}
\]


\[
X^TX = \begin{pmatrix} 3 & 6 \\ 6 & 14 \end{pmatrix}, \qquad X^T\vec{y} = \begin{pmatrix} 2+3+5 \\ 1\cdot 2 + 2\cdot 3 + 3\cdot 5 \end{pmatrix} = \begin{pmatrix} 10 \\ 23 \end{pmatrix}
\]


b) Die analytische Normalgleichung liefert das \textbf{unbeschränkte} Optimum. Sobald NB vorliegen ($b_i \ge 0$, $\sum b_i = 4$), kann diese Lösung außerhalb des zulässigen Bereichs liegen, also ist sie ungültig. Genau diese Situation definiert laut Kapitel das CLS-Problem, und wir müssen es als QO mit NB lösen.


c) Aus $\|\vec{y} - X\vec{b}\|^2 = \vec{y}^T\vec{y} - 2\vec{y}^TX\vec{b} + \vec{b}^TX^TX\vec{b}$ folgt durch Weglassen der Konstanten und Multiplikation mit $\tfrac{1}{2}$:


\[
\min_{\vec{b}}\ \tfrac{1}{2}\vec{b}^T\underbrace{(X^TX)}_{Q}\vec{b}\ +\ \underbrace{(-X^T\vec{y})^T}_{\vec{c}^{\,T}}\vec{b}
\]


\[
Q = \begin{pmatrix} 3 & 6 \\ 6 & 14 \end{pmatrix},\qquad \vec{c} = \begin{pmatrix} -10 \\ -23 \end{pmatrix}
\]


Restriktionen:

\begin{itemize}
  \item Gleichungs-NB: $A_{eq}\vec{b} = b_{eq}$ mit $A_{eq} = (1,\ 1)$, $b_{eq} = 4$.
  \item Box-NB: $b_l = (0,0)^T \le \vec{b}$ (kein oberer Wert nötig).
\end{itemize}

$Q$ hat $\det = 42 - 36 = 6 > 0$ und $Q_{11} = 3 > 0$, also positiv definit → konvexes QO-Problem mit eindeutiger Lösung.



\noindent\textcolor{rulegray}{\rule{\linewidth}{0.4pt}}


\paragraph{Aufgabe 3 — KKT-Bedingungen anwenden \textit{(20 Punkte)}}\mbox{}\\[2pt]

Gegeben:


\[
\min f(x,y) = x^2 + 2y^2,\qquad \text{u.B.}\quad x+y = 3,\ \ x \ge 0,\ \ y \ge 0
\]


a) Bringe die Ungleichungs-NB in die Normalform $g_j(\vec{x}) \le 0$ und schreibe die Lagrange-Funktion $L(x,y,a,b_1,b_2)$ an.

b) Stelle die vier KKT-Blöcke vollständig auf.

c) Bestimme den KKT-Punkt. Mache dazu eine begründete Fallunterscheidung über die aktiven Ungleichungs-NB und prüfe alle KKT-Bedingungen am Ergebnis.

d) Argumentiere mit der Bedingung 2. Ordnung, ob es sich um ein lokales Minimum handelt.


\textbf{Lösung:}


a) Normalform: $g_1(\vec{x}) = -x \le 0$, $g_2(\vec{x}) = -y \le 0$, $h_1(\vec{x}) = x+y-3 = 0$.


\[
L = x^2 + 2y^2 + a(x+y-3) + b_1(-x) + b_2(-y)
\]


b) KKT-Blöcke:


1) \textbf{Stationarität}:

\[
\frac{\partial L}{\partial x} = 2x + a - b_1 = 0,\qquad \frac{\partial L}{\partial y} = 4y + a - b_2 = 0
\]

2) \textbf{Primal Feasibility}: $x+y-3 = 0$, $-x \le 0$, $-y \le 0$.

3) \textbf{Dual Feasibility}: $b_1 \ge 0$, $b_2 \ge 0$.

4) \textbf{Komplementarität}: $b_1\cdot(-x) = 0$, $b_2\cdot(-y) = 0$.


c) \textbf{Fallunterscheidung}:


\textit{Annahme: beide Ungleichungs-NB inaktiv}, also $b_1 = b_2 = 0$ (Komplementarität, falls $x>0, y>0$).


Dann liefert die Stationarität:

\[
2x + a = 0\ \Rightarrow\ x = -a/2,\qquad 4y + a = 0\ \Rightarrow\ y = -a/4
\]


Einsetzen in die Gleichungs-NB:

\[
x+y = -\tfrac{a}{2} - \tfrac{a}{4} = -\tfrac{3a}{4} = 3 \Rightarrow a = -4
\]


\[
x^* = 2,\quad y^* = 1
\]


Prüfung aller KKT:

\begin{itemize}
  \item Stationarität: $2\cdot 2 + (-4) = 0$ ✓, $4\cdot 1 + (-4) = 0$ ✓
  \item Primal: $2+1=3$ ✓, $-2 \le 0$ ✓, $-1 \le 0$ ✓
  \item Dual: $b_1 = b_2 = 0 \ge 0$ ✓
  \item Komplementarität: $0\cdot(-2) = 0$ ✓, $0\cdot(-1) = 0$ ✓
\end{itemize}

KKT-Punkt: $(x^*, y^*, a, b_1, b_2) = (2,\ 1,\ -4,\ 0,\ 0)$, $f^* = 4 + 2 = 6$.


(Andere Fälle: aktiv $g_1$ ($x=0$) liefert $y=3$ → $b_1 = a$, aus $4y+a=0$ folgt $a=-12$, $b_1=-12 < 0$ → verletzt Dual Feasibility. Analog für $g_2$. Damit ist obiger Punkt der einzige KKT-Punkt.)


d) Hesse von $L$ bzgl. $\vec{x}$:

\[
\nabla^2_{\vec{x}} L = \begin{pmatrix} 2 & 0 \\ 0 & 4 \end{pmatrix}
\]

positiv definit (Hauptminoren $2>0$, $8>0$). Damit ist das hinreichende Kriterium 2. Ordnung erfüllt → $(2,1)$ ist ein \textbf{lokales Minimum} (sogar global, da $f$ konvex und der zulässige Bereich konvex ist).



\noindent\textcolor{rulegray}{\rule{\linewidth}{0.4pt}}


\paragraph{Aufgabe 4 — Eine Iteration SQO \textit{(20 Punkte)}}\mbox{}\\[2pt]

Gegeben sei das NLO


\[
\min f(x_1,x_2) = x_1^2 + x_2^2 + x_1 x_2,\qquad \text{u.B.}\ x_1 + x_2 \ge 2
\]


mit Startpunkt $\vec{x}_0 = (3, 0)^T$.


a) Schreibe die Ungleichungs-NB als $g(\vec{x}) \le 0$. Berechne $\nabla f(\vec{x}_0)$, $\nabla g(\vec{x}_0)$, $g(\vec{x}_0)$.

b) Stelle das QO-Sub-Problem nach Kapitel-Definition auf (mit $H$ = exakte Hesse von $L$ bzgl. $\vec{x}$ am Punkt $\vec{x}_0$ — begründe, dass dies hier zulässig ist).

c) Löse das Sub-Problem über KKT (Sub-Problem hat eigene KKT-Bedingungen!) und bestimme die Suchrichtung $\vec{d}$.

d) Führe einen vollständigen Iterationsschritt mit Schrittweite $c=1$ durch. Vergleiche $f(\vec{x}_0)$ und $f(\vec{x}_1)$ und prüfe, ob $\vec{x}_1$ zulässig ist.


\textbf{Lösung:}


a) Normalform: $g(\vec{x}) = 2 - x_1 - x_2 \le 0$.


\[
\nabla f(\vec{x}) = \begin{pmatrix} 2x_1 + x_2 \\ 2x_2 + x_1 \end{pmatrix},\quad \nabla f(\vec{x}_0) = \begin{pmatrix} 6 \\ 3 \end{pmatrix}
\]


\[
\nabla g(\vec{x}_0) = \begin{pmatrix} -1 \\ -1 \end{pmatrix},\qquad g(\vec{x}_0) = 2 - 3 - 0 = -1\ (<0,\ \text{NB inaktiv})
\]


b) Hesse:

\[
\nabla^2 f = \begin{pmatrix} 2 & 1 \\ 1 & 2 \end{pmatrix},\qquad \nabla^2 g = 0
\]

Da $g$ linear ist, ist $\nabla^2_{\vec{x}} L = \nabla^2 f + b\cdot\nabla^2 g = \nabla^2 f$ unabhängig von $b$. Damit ist $H = \begin{pmatrix} 2 & 1 \\ 1 & 2 \end{pmatrix}$ exakt und positiv definit (Hauptminoren $2>0$, $3>0$); BFGS wäre nicht nötig.


Sub-Problem ($\vec{d} = (d_1,d_2)^T$):

\[
\min_{\vec{d}}\ f(\vec{x}_0) + \nabla f^T\vec{d} + \tfrac{1}{2}\vec{d}^T H \vec{d}
\]

\[
\text{u.B.}\quad g(\vec{x}_0) + \nabla g^T\vec{d} \le 0\ \Leftrightarrow\ -1 - d_1 - d_2 \le 0\ \Leftrightarrow\ d_1 + d_2 \ge -1
\]


Konstanten weglassen:

\[
\min\ \tfrac{1}{2}(2d_1^2 + 2d_2^2 + 2d_1 d_2) + 6d_1 + 3d_2 \quad\text{s.t.}\ -1 - d_1 - d_2 \le 0
\]


c) KKT mit Multiplikator $b \ge 0$:


\[
2d_1 + d_2 + 6 - b = 0,\quad d_1 + 2d_2 + 3 - b = 0,\quad b\cdot(-1-d_1-d_2) = 0
\]


\textbf{Versuch $b=0$} (NB inaktiv):

Aus den ersten beiden Gleichungen: $2d_1+d_2=-6$, $d_1+2d_2=-3$. Subtraktion liefert $d_1 - d_2 = -3$, eingesetzt: $d_2 = 0$, $d_1 = -3$. Aber $-1 - (-3) - 0 = 2 > 0$ — NB \textbf{verletzt}.


\textbf{Damit $b>0$, NB aktiv}: $d_1 + d_2 = -1$, also $d_2 = -1 - d_1$.


\[
2d_1 + (-1-d_1) + 6 - b = 0\ \Rightarrow\ d_1 + 5 = b
\]

\[
d_1 + 2(-1-d_1) + 3 - b = 0\ \Rightarrow\ -d_1 + 1 = b
\]


Gleichsetzen: $d_1 + 5 = -d_1 + 1\ \Rightarrow\ d_1 = -2$, $d_2 = 1$, $b = 3 > 0$ ✓.


Suchrichtung: $\vec{d} = (-2,\ 1)^T$.


d) Mit $c=1$:

\[
\vec{x}_1 = \vec{x}_0 + 1\cdot \vec{d} = (3,0)^T + (-2,1)^T = (1,1)^T
\]


Zulässigkeit: $x_1 + x_2 = 2 \ge 2$ ✓ (NB jetzt aktiv).


Funktionswerte: $f(\vec{x}_0) = 9 + 0 + 0 = 9$, $f(\vec{x}_1) = 1 + 1 + 1 = 3$ → echter Abstieg um 6.


(Anmerkung: $(1,1)$ ist hier sogar bereits das globale Optimum des ursprünglichen NLO; in einer zweiten Iteration würde $\vec{d} \approx \vec{0}$ und der Algorithmus konvergieren.)



\noindent\textcolor{rulegray}{\rule{\linewidth}{0.4pt}}


\paragraph{Aufgabe 5 — Fehler im SQO-Pseudocode \textit{(15 Punkte)}}\mbox{}\\[2pt]

Ein Studierender hat die folgende Implementierung des SQO-Algorithmus geschrieben:


\begin{lstlisting}
INPUT: NLO mit f, h, g, Startpunkt x_0
1:  i ← 0
2:  H ← I              // I = Einheitsmatrix
3:  while i < 100 do
4:      grad_f ← ∇f(x_i)
5:      grad_h ← ∇h(x_i);  grad_g ← ∇g(x_i)
6:      löse QO-Sub-Problem mit grad_f, grad_h, grad_g, H → d
7:      x_{i+1} ← x_i + d
8:      i ← i + 1
9:  end while
10: return x_i
\end{lstlisting}


a) Identifiziere mindestens \textbf{drei} inhaltliche Fehler/Mängel gegenüber dem Kapitel und erkläre jeweils, welche praktische Konsequenz der Fehler hat.

b) Formuliere den Algorithmus in korrigierter Form (Pseudocode oder Stichpunkte).

c) Begründe, warum das in (a) gewählte Verfahren zur Berechnung von $H$ insbesondere bei nicht-linearen Problemen wichtig ist.


\textbf{Lösung:}


a) Mängel:


\begin{enumerate}
  \item \textbf{Fixe Hesse $H = I$}: Im Kapitel ist $H$ eine Approximation der Hesse von $L$ bzgl. $\vec{x}$ (typisch BFGS-Update). Mit $H = I$ wird das Sub-Problem zu einer skalierten Gradientenrichtung — die quadratische Krümmung von $f$ und das Verhalten der NB-Multiplikatoren werden vollständig ignoriert. Konsequenz: extrem langsame oder gar fehlende Konvergenz, kein superlinearer Charakter mehr.
  \item \textbf{Fehlende Liniensuche}: Das Kapitel verlangt eine Schrittgröße $c > 0$ aus einer Liniensuche und $\vec{x}_{i+1} = \vec{x}_i + c\cdot\vec{d}$. Im Code wird stets $c=1$ verwendet. Konsequenz: Schritt kann zu groß sein → Zielfunktion steigt, NB werden grob verletzt, Algorithmus divergiert oder oszilliert.
  \item \textbf{Falsches Abbruchkriterium}: Es wird nur die Iterationszahl geprüft, kein Konvergenzkriterium (z. B. $\|\vec{d}\| < \varepsilon$, $\|\nabla L\| < \varepsilon$, KKT-Residuum). Konsequenz: Entweder zu früh abgebrochen (kein lokales Optimum) oder unnötig lange weitergerechnet, ohne Garantie über die Qualität von $\vec{x}_i$.
\end{enumerate}

(Optional weitere Punkte: keine Multiplikatoren-Übernahme aus dem Sub-Problem zur $H$-Update / Liniensuche; keine Prüfung der Lösbarkeit des Sub-Problems; keine Zulässigkeitsprüfung von $\vec{x}_0$.)


b) Korrigierte Fassung:


\begin{lstlisting}
INPUT: NLO (f, h, g), Startpunkt x_0, Toleranz ε > 0, max_iter
1:  i ← 0;  H ← I  (Initialisierung; wird via BFGS aktualisiert)
2:  repeat
3:      grad_f ← ∇f(x_i);  grad_h ← ∇h(x_i);  grad_g ← ∇g(x_i)
4:      löse QO-Sub-Problem → Suchrichtung d, Multiplikatoren a, b
5:      bestimme Schrittgröße c > 0 via Liniensuche (z. B. Armijo / Merit-Funktion)
6:      x_{i+1} ← x_i + c·d
7:      aktualisiere H via BFGS aus (x_i, x_{i+1}, ∇L)
8:      i ← i + 1
9:  until ‖d‖ < ε  ODER  KKT-Residuum < ε  ODER  i ≥ max_iter
10: return x_i, a, b
\end{lstlisting}


c) Bei nicht-linearen Problemen variiert die Krümmung von $f$ und $L$ über den Iterationsverlauf. Eine Approximation der Hesse von $L$ (statt eines reinen Gradientenschritts) sorgt dafür, dass das Sub-Problem die \textit{lokale} Krümmung berücksichtigt — dadurch wird die Suchrichtung lokal nahezu Newton-artig und SQO erreicht typischerweise superlineare Konvergenz, was bei $H=I$ verlorengeht. BFGS wird dabei der direkten Hessen-Berechnung vorgezogen, weil zweite Ableitungen oft teuer oder nicht analytisch verfügbar sind.





% ──────────────────────────────────────────────────────────────

\end{document}
